\chapter{The Derived Category}

\section{Introduction}
Let \(\mathcal{A}\) and \(\mathcal{B}\) be abelian categories, and let \(F\colon \mathcal{A}\to\mathcal{B}\) be a functor. Our purpose is to define, functorially for each \(X\in\operatorname{Ob}\mathcal{A}\), a complex \(\mathbb{R}F(X)\) whose cohomology groups are the right derived functors of \(F\) acting on \(X\), namely \(R^iF(X)\). In general, we will define \(\mathbb{R}F(X^\bullet)\) for any complex \(X^\bullet\) of objects of \(\mathcal{A}\).

More precisely, \(\mathbb{R}F\) will be a functor from the derived category \(\mathcal{D}(\mathcal{A})\) of \(\mathcal{A}\) to the derived category \(\mathcal{D}(\mathcal{B})\) of \(\mathcal{B}\).

The derived category \(\mathcal{D}(\mathcal{A})\) is obtained as follows: one first considers the category \(\mathcal{K}(\mathcal{A})\), whose objects are complexes of elements of \(\mathcal{A}\), and whose morphisms are homotopy equivalence classes of morphisms of complexes. The category \(\mathcal{D}(\mathcal{A})\) is obtained by ``localizing'' \(\mathcal{K}(\mathcal{A})\) so that every morphism in \(\mathcal{K}(\mathcal{A})\) which induces an isomorphism on cohomology becomes an isomorphism in \(\mathcal{D}(\mathcal{A})\). This process of localization will be explained in general below.

Although \(\mathcal{A}\) and \(\mathcal{B}\) above are assumed to be abelian categories, the categories \(\mathcal{K}(\mathcal{A})\), \(\mathcal{D}(\mathcal{A})\), etc., are in general not abelian. However, they can be given a structure which carries enough information for our purposes, namely a structure of triangulated category. Thus we will be led to study triangulated categories.

\section{Triangulated Categories}
\subsection{Definition}
A triangulated category is an additive category \(\mathcal{C}\), together with
\begin{enumerate}
	\item[a)] an automorphism \(T\colon \mathcal{C}\to\mathcal{C}\) called the \emph{translation functor}, and
	\item[b)] a collection of sextuples \((X,Y,Z,u,v,w)\), called the \emph{triangles} of \(\mathcal{C}\),
\end{enumerate}
where in each sextuple \(X,Y,Z\in\operatorname{Ob}\mathcal{C}\) and
\[
u\colon X\to Y,\quad v\colon Y\to Z,\quad w\colon Z\to T(X)
\]
are morphisms. A triangle is usually written
\[
X\xrightarrow{u} Y\xrightarrow{v} Z\xrightarrow{w} T(X).
\]

A morphism of triangles
\[
(X,Y,Z,u,v,w)\to(X',Y',Z',u',v',w')
\]
is a commutative diagram
\[
\begin{CD}
	X @>u>> Y @>v>> Z @>w>> T(X) \\
	@VfVV @VgVV @VhVV @VT(f)VV \\
	X' @>u'>> Y' @>v'>> Z' @>w'>> T(X')
\end{CD}
\]

This data is subject to the following axioms:

\begin{enumerate}
	\item[\textbf{(TR1)}]
	\begin{enumerate}
		\item Every sextuple isomorphic to a triangle is a triangle.
		\item Every morphism \(u\colon X\to Y\) can be embedded in a triangle.
		\item The sextuple \((X,X,0,\operatorname{id}_X,0,0)\) is a triangle.
	\end{enumerate}
	
	\item[\textbf{(TR2)}]
	\((X,Y,Z,u,v,w)\) is a triangle if and only if \((Y,Z,T(X),v,w,-T(u))\) is a triangle.
	
	\item[\textbf{(TR3)}]
	Given two triangles \((X,Y,Z,u,v,w)\) and \((X',Y',Z',u',v',w')\), and morphisms \(f\colon X\to X'\), \(g\colon Y\to Y'\) such that \(u'f=gu\), there exists a morphism \(h\colon Z\to Z'\) (not necessarily unique) so that \((f,g,h)\) is a morphism of triangles.
	
	\item[\textbf{(TR4)}] (The octahedral axiom.)\\
	\begin{center}
	\begin{tikzcd}
		& Y' \arrow[dr, dashed, "g"] & \\
		Z' \arrow[ur, dashed, "f"] \arrow[d] & & X' \arrow[ll, "T(j)i"] \arrow[d] \\
		X \arrow[rr, "vu"] \arrow[dr, "u"'] \arrow[uur, "j" near start] & & Z \arrow[u] \arrow[dl, "v"'] \arrow[uul, "i" near start] \\
		& Y \arrow[uul] \arrow[uur] & 
	\end{tikzcd}
	\end{center}
	
	Suppose given triangles
	\[
	\begin{gathered}
		(X,Y,Z',u,j,\cdot),\\
		(Y,Z,X',v,\cdot,i),\\
		(X,Z,Y',vu,\cdot,\cdot)
	\end{gathered}
	\]
	such that
	\[
	(Z',Y',X',f,g,T(j)i)
	\]
	is a triangle. Then there exist morphisms \(f\colon Z'\to Y'\) and \(g\colon Y'\to X'\) such that the two other faces of the octahedron with \(f,g\) as edges are commutative diagrams.
\end{enumerate}

\begin{definition}
	An additive functor \(F\colon \mathcal{C}\to\mathcal{C}'\) from one triangulated category to another is called a \emph{(covariant) \(\partial\)-functor} if it commutes with the translation functor and takes triangles into triangles. A contravariant \(\partial\)-functor takes triangles into triangles with the arrows reversed and sends the translation functor into its inverse.
\end{definition}

\begin{definition}
	An additive functor \(H\colon \mathcal{C}\to\mathcal{A}\) from a triangulated category to an abelian category is called a \emph{covariant cohomological functor} if whenever \((X,Y,Z,u,v,w)\) is a triangle, the long sequence
	\[
	\cdots\to H(T^iX)\to H(T^iY)\to H(T^iZ)\to H(T^{i+1}X)\to\cdots
	\]
	is exact (the morphisms being \(H(T^iu)\), etc.). If \(H\) is a cohomological functor, we often write \(H^i(X)\) for \(H(T^iX)\), \(i\in\mathbb{Z}\). One defines a contravariant cohomological functor by reversing the arrows.
\end{definition}

\begin{proposition}\label{prop:1.1}
	\begin{enumerate}
		\item[a)] The composition of any two consecutive morphisms in a triangle is zero.
		\item[b)] If \(\mathcal{C}\) is a triangulated category and \(M\) an object of \(\mathcal{C}\), then \(\operatorname{Hom}_{\mathcal{C}}(M,\cdot)\) and \(\operatorname{Hom}_{\mathcal{C}}(\cdot,M)\) are cohomological functors on \(\mathcal{C}\).
		\item[c)] If in the situation of (TR3) \(f\) and \(g\) are isomorphisms, then \(h\) is also an isomorphism.
	\end{enumerate}
\end{proposition}

\begin{proof}
	a) Let \((X,Y,Z,u,v,w)\) be a triangle. By (TR2) it is sufficient to show that \(vu=0\). By (TR1), \((Z,Z,0,\operatorname{id}_Z,0,0)\) is a triangle. Also by (TR2), \((Y,Z,T(X),v,w,-T(u))\) is a triangle. We apply (TR3) to the maps \(v\colon Y\to Z\) and \(\operatorname{id}_Z\colon Z\to Z\), giving a morphism of triangles. It follows that \(T(v)(-T(u))=0\), or since \(T\) is an automorphism, \(vu=0\).
	
	b) Let \((X,Y,Z,u,v,w)\) be a triangle. To show \(\operatorname{Hom}_{\mathcal{C}}(M,\cdot)\) is a cohomological functor, it suffices by (TR2) to show the sequence
	\[
	\operatorname{Hom}_{\mathcal{C}}(M,X)\to\operatorname{Hom}_{\mathcal{C}}(M,Y)\to\operatorname{Hom}_{\mathcal{C}}(M,Z)
	\]
	is exact. By a) we know the composition is zero. Now suppose given \(g\in\operatorname{Hom}_{\mathcal{C}}(M,Y)\) such that \(vg=0\). We apply (TR3) to the triangles \((M,M,0,\operatorname{id}_M,0,0)\) and \((X,Y,Z,u,v,w)\) and the map \(g\colon M\to Y\) to conclude there exists \(f\colon M\to X\) such that \(uf=g\).
	
	A similar proof shows that \(\operatorname{Hom}_{\mathcal{C}}(\cdot,M)\) is a contravariant cohomological functor.
	
	c) In the situation of (TR3), suppose \(f,g\) are isomorphisms. Apply the cohomological functor \(\operatorname{Hom}_{\mathcal{C}}(Z',\cdot)\) to obtain an exact commutative diagram where \(f_*,g_*,T(f)_*,T(g)_*\) are isomorphisms of abelian groups. By the five-lemma, \(h_*\) is an isomorphism. Hence there exists \(\varphi\in\operatorname{Hom}_{\mathcal{C}}(Z',Z)\) such that \(h\varphi=\operatorname{id}_Z\). Using \(\operatorname{Hom}_{\mathcal{C}}(\cdot,Z)\) similarly gives \(\psi\in\operatorname{Hom}_{\mathcal{C}}(Z,Z')\) such that \(\psi h=\operatorname{id}_{Z'}\). Thus \(h\) is an isomorphism.
\end{proof}

\section{\(\mathcal{K}(\mathcal{A})\) is triangulated}
Let \(\mathcal{A}\) be an abelian category. A \emph{complex} of objects of \(\mathcal{A}\) is a collection \(X^\bullet=(X^n)_{n\in\mathbb{Z}}\) of objects of \(\mathcal{A}\) together with maps \(d^n\colon X^n\to X^{n+1}\) such that \(d^{n+1}d^n=0\) for all \(n\in\mathbb{Z}\).

A morphism \(f\) of complexes \(X^\bullet\to Y^\bullet\) is a collection of maps \(f^n\colon X^n\to Y^n\) which commute with differentials:
\[
f^{n+1}d_X^n = d_Y^n f^n
\]
for all \(n\).

Two maps \(f,g\colon X^\bullet\to Y^\bullet\) are \emph{homotopic} if there exists a collection of maps \(k=(k^n)\), \(k^n\colon X^n\to X^{n-1}\) (not necessarily commuting with \(d\)) such that
\[
f^n - g^n = d_Y^{n-1}k^n + k^{n+1}d_X^n
\]
for all \(n\). Homotopy is an equivalence relation, and compositions of homotopic maps are homotopic.

We define \(\mathcal{K}(\mathcal{A})\) to be the category whose objects are complexes of objects of \(\mathcal{A}\), and whose morphisms are homotopy equivalence classes of morphisms of complexes.

A complex \(X^\bullet\) is \emph{bounded below} if \(X^n=0\) for \(n\ll0\). We denote by \(\mathcal{K}^+(\mathcal{A})\) the full subcategory of \(\mathcal{K}(\mathcal{A})\) of complexes bounded below. Similarly we define \(\mathcal{K}^-(\mathcal{A})\) (bounded above) and \(\mathcal{K}^b(\mathcal{A})\) (bounded on both sides).

We now define a triangulated structure on \(\mathcal{K}(\mathcal{A})\) (resp.\ \(\mathcal{K}^+(\mathcal{A})\), etc.). Let \(T\) be the left shift functor with sign change:
\[
T(X^\bullet)^p = X^{p+1},\qquad d_{T(X)} = -d_X.
\]
We often write \(X^\bullet[1]\) instead of \(T(X^\bullet)\), and \(X^\bullet[n]\) instead of \(T^n(X^\bullet)\).

If \(u\colon X^\bullet\to Y^\bullet\) is any morphism, consider its \emph{mapping cone} \(Z^\bullet\), defined as the complex
\[
Z^\bullet = T(X^\bullet)\oplus Y^\bullet
\]
with differential
\[
d_Z = \begin{pmatrix} T(d_X) & T(u) \\ 0 & d_Y \end{pmatrix}.
\]
There are natural morphisms
\[
v\colon Y^\bullet\to Z^\bullet,\qquad w\colon Z^\bullet\to T(X^\bullet).
\]

We define a triangle in \(\mathcal{K}(\mathcal{A})\) to be any sextuple isomorphic to one obtained from a morphism \(u\) by this construction.

All axioms (TR1)--(TR4) are easily verified, once one notes that the mapping cone of \(\operatorname{id}_{X^\bullet}\) is homotopic to \(0\). Indeed, the matrix
\[
k = \begin{pmatrix} 0 & 0 \\ \operatorname{id}_X & 0 \end{pmatrix}
\]
gives a contraction homotopy.

\begin{definition}
	We define \(H\) to be the functor from \(\mathcal{K}(\mathcal{A})\) to \(\mathcal{A}\) which sends a complex \(X^\bullet\) to its \(0^\text{th}\) cohomology:
	\[
	H(X^\bullet) = \ker d^0 \big/ \operatorname{im} d^{-1}.
	\]
	This is indeed a functor, since homotopic maps induce the same map on cohomology. We write \(H^i = H\circ T^i\) for any \(i\in\mathbb{Z}\).
\end{definition}

Observe that \(H\) is a cohomological functor from \(\mathcal{K}(\mathcal{A})\) to \(\mathcal{A}\). Indeed, it suffices to check the long exact sequence for triangles coming from mapping cones, which is straightforward.

\section{Localization of categories}

\begin{definition}
Let \(\mathcal{C}\) be a category. A collection \(S\) of morphisms of \(\mathcal{C}\) is called a \emph{multiplicative system} if it satisfies the following axioms (FR1)--(FR3):
\begin{enumerate}
\item[(FR1)] If \(f,g\in S\) and the composite \(fg\) exists, then \(fg\in S\). For any \(X\in\operatorname{Ob}\mathcal{C}\), \(\operatorname{id}_X\in S\).

\item[(FR2)] Any diagram
\[
X\xrightarrow{u} Y \xleftarrow{s} Z,\quad s\in S
\]
can be completed to a commutative diagram
\[
\begin{CD}
W @>v>> Z\\
@VtVV @VsVV\\
X @>u>> Y
\end{CD}
\]
with \(t\in S\). The analogous statement holds with all arrows reversed.

\item[(FR3)] Let \(f,g\colon X\to Y\) be morphisms in \(\mathcal{C}\). The following are equivalent:
\begin{enumerate}
\item[(i)] There exists \(s\colon Y\to Y'\) in \(S\) such that \(sf = sg\);
\item[(ii)] There exists \(t\colon X'\to X\) in \(S\) such that \(ft = gt\).
\end{enumerate}
\end{enumerate}
\end{definition}

\begin{definition}
Let \(\mathcal{C}\) be a category and \(S\) a multiplicative system of morphisms in \(\mathcal{C}\). The \emph{localization of \(\mathcal{C}\) with respect to \(S\)} is a category \(\mathcal{C}_S\) together with a functor \(Q\colon\mathcal{C}\to\mathcal{C}_S\) such that:
\begin{enumerate}
\item[(a)] \(Q(s)\) is an isomorphism for every \(s\in S\);
\item[(b)] Any functor \(F\colon\mathcal{C}\to\mathcal{D}\) sending each \(s\in S\) to an isomorphism factors uniquely through \(Q\).
\end{enumerate}
\end{definition}

\begin{proposition}\label{prop:1.3.1}
Let \(\mathcal{C}\) be a category and \(S\) a multiplicative system in \(\mathcal{C}\). The localization \(\mathcal{C}_S\) may be constructed as follows:
\(\operatorname{Ob}\mathcal{C}_S = \operatorname{Ob}\mathcal{C}\), and for all \(X,Y\in\operatorname{Ob}\mathcal{C}\),
\[
\operatorname{Hom}_{\mathcal{C}_S}(X,Y) = \varinjlim_{I_X} \operatorname{Hom}_{\mathcal{C}}(X',Y),
\]
where \(I_X\) is the category whose objects are morphisms \(s\colon X'\to X\) in \(S\), and morphisms are commutative triangles over \(X\).

Furthermore, if \(\mathcal{C}\) is additive, then \(\mathcal{C}_S\) is also additive.
\end{proposition}

\begin{proof}
Using (FR1), (FR2), (FR3), the index category \(I_X\) satisfies the inductive limit axioms (L1, L2, L3) of [GT, Ch.\ I, \S1], so the direct limit is well-behaved.

A morphism \(X\to Y\) in \(\mathcal{C}_S\) is represented by a right fraction \(X\xleftarrow{s}X'\xrightarrow{f}Y\) with \(s\in S\). Two such fractions define the same morphism iff there exist \(u\colon X'''\to X'\), \(g\colon X'''\to X''\) in \(S\) making the obvious diagram commute and equalizing the representatives.

Composition of fractions is defined via (FR2): given fractions \(X\xleftarrow{s}X'\xrightarrow{a}Y\) and \(Y\xleftarrow{t}Y'\xrightarrow{b}Z\), complete the diagram to get a common lift and form the composite morphism \(X\to Z\). The result is independent of representative and lift.

The localization functor \(Q\colon\mathcal{C}\to\mathcal{C}_S\) satisfies the universal property, and additivity of \(\mathcal{C}\) passes to \(\mathcal{C}_S\) since the limit preserves abelian group structure on Hom-sets.
\end{proof}

\begin{definition}
Let \(\mathcal{C}\) be a triangulated category and \(S\) a multiplicative system in \(\mathcal{C}\). We say \(S\) is \emph{compatible with the triangulation} if:
\begin{enumerate}
\item[(FR4)] \(s\in S \implies T(s)\in S\), where \(T\) is the translation functor;
\item[(FR5)] Axiom (TR3) holds with the extra condition: if \(f,g\in S\) are the given morphisms of triangles, then one may choose the filler \(h\in S\).
\end{enumerate}
\end{definition}

\begin{proposition}\label{prop:1.3.2}
If \(\mathcal{C}\) is triangulated and \(S\) a triangulation-compatible multiplicative system, then \(\mathcal{C}_S\) admits a \emph{unique} triangulated structure such that \(Q\colon\mathcal{C}\to\mathcal{C}_S\) is a triangle functor, and \(Q\) satisfies the universal property for triangle functors into other triangulated categories.
\end{proposition}

One may also compute morphism sets as
\[
\operatorname{Hom}_{\mathcal{C}_S}(X,Y) = \varinjlim_{J_Y}\operatorname{Hom}_{\mathcal{C}}(X,Y'),
\]
where \(J_Y\) has objects \(s\colon Y\to Y'\in S\) and morphisms commutative diagrams under \(Y\).

\begin{proposition}\label{prop:1.3.3}
Let \(\mathcal{C}\) be a category, \(S\) a multiplicative system in \(\mathcal{C}\), and \(\mathcal{D}\subseteq\mathcal{C}\) a full subcategory such that \(S\cap\mathcal{D}\) is a multiplicative system in \(\mathcal{D}\). Assume one of the following:
\begin{enumerate}
\item[(i)] For every \(s\colon X'\to X\) in \(S\) with \(X\in\operatorname{Ob}\mathcal{D}\), there exists \(f\colon X''\to X'\) with \(X''\in\operatorname{Ob}\mathcal{D}\) and \(sf\in S\);
\item[(ii)] The same with all arrows reversed.
\end{enumerate}
Then the natural functor
\[
\mathcal{D}_{S\cap\mathcal{D}} \to \mathcal{C}_S
\]
is fully faithful, so \(\mathcal{D}_{S\cap\mathcal{D}}\) may be identified with a full subcategory of \(\mathcal{C}_S\).
\end{proposition}

\begin{proposition}\label{prop:1.3.4}
Let \(\mathcal{C}\) be a category, \(S\) a multiplicative system, \(Q\colon\mathcal{C}\to\mathcal{C}_S\) the localization functor. Let \(\mathcal{D}\) be any category, and \(F,G\colon\mathcal{C}_S\to\mathcal{D}\) functors. Then the natural restriction map
\[
\alpha\colon \operatorname{Hom}(F,G) \to \operatorname{Hom}(FQ,GQ)
\]
between functor morphism sets is bijective.
\end{proposition}

\begin{proof}
A functor morphism \(F\to G\) assigns to each object \(X\in\mathcal{C}_S\) a morphism \(F(X)\to G(X)\) making all morphism squares commute.

Injectivity of \(\alpha\) is clear since \(\operatorname{Ob}\mathcal{C}_S=\operatorname{Ob}\mathcal{C}\). For surjectivity, take a morphism \(FQ\to GQ\). Since every morphism in \(\mathcal{C}_S\) is a fraction built from morphisms of \(\mathcal{C}\) and elements of \(S\), and \(F(s),G(s)\) are isomorphisms for \(s\in S\), the given functor morphism on \(FQ,GQ\) extends uniquely to a functor morphism on \(F,G\) over \(\mathcal{C}_S\).
\end{proof}

\section{Quasi-isomorphisms and the derived category}

Let \(\mathcal{A}\) be an abelian category, and let \(\KK(\mathcal{A})\) be the triangulated category of complexes up to homotopy introduced in \S2. A morphism \(f\colon X^\bullet\to Y^\bullet\) in \(\KK(\mathcal{A})\) is called a \emph{quasi-isomorphism} if it induces an isomorphism on all cohomology groups. Denote by \(\operatorname{Qis}\) the collection of all quasi-isomorphisms in \(\KK(\mathcal{A})\).

\begin{proposition}\label{prop:1.4.2}
Let \(\mathcal{C}\) be a triangulated category, \(\mathcal{A}\) an abelian category, and \(H\colon\mathcal{C}\to\mathcal{A}\) a cohomological functor. Let \(S\) be the class of morphisms \(s\in\mathcal{C}\) such that \(H(T^i(s))\) is an isomorphism for all \(i\in\mathbb{Z}\). Then \(S\) is a multiplicative system in \(\mathcal{C}\) and is compatible with the triangulation.
\end{proposition}

\begin{proof}
We verify axioms (FR1)--(FR5). Axioms (FR1) and (FR4) are trivial. Axiom (FR5) follows from the long exact sequence of a cohomological functor and the five-lemma.

To prove (FR2), consider a diagram
\[
X\xrightarrow{u} Y \xleftarrow{s} Z,\qquad s\in S.
\]
Complete \(s\) to a triangle \((Z,Y,N,s,f,g)\). Complete \(fu\) to a triangle \((W,X,N,t,fu,h)\). Since \((u,\operatorname{id}_N)\) maps two sides of the second triangle to the first, there exists a morphism \(v\colon W\to Z\) giving a morphism of triangles, satisfying \(sv = ut\).

It remains to show \(t\in S\). Since \(s\in S\), the long exact sequence of the first triangle implies \(H(T^i(N))=0\) for all \(i\in\mathbb{Z}\). Substituting into the long exact sequence of the second triangle yields that \(H(T^i(t))\) is an isomorphism for all \(i\in\mathbb{Z}\), so \(t\in S\). The dual statement of (FR2) is proved similarly.

For (FR3), consider \(f,g\colon X\to Y\). It suffices to prove equivalence of:
\begin{enumerate}
\item[(i)] \(\exists\;s\colon Y\to Y'\in S\) with \(sf=sg\);
\item[(ii)] \(\exists\;t\colon X'\to X\in S\) with \(ft=gt\).
\end{enumerate}
It is enough to treat the case \(g=0\), so we need:
(i) \(\exists\,s\colon Y\to Y'\in S,\ sf=0\)
\(\iff\)
(ii) \(\exists\,t\colon X'\to X\in S,\ ft=0\).

Assume (i). By (TR1),(TR2), form a triangle \((Z,Y,Y',v,s,u)\). Since \(sf=0\), Proposition 1.1(b) gives \(g\colon X\to Z\) with \(f=vg\). Form a triangle \((X',X,Z,t,g,w)\). Again by Proposition 1.1(b), \(ft=0\). Since \(s\in S\), \(H(T^i(Z))=0\) for all \(i\), which implies \(t\in S\). The converse (ii)\(\Rightarrow\)(i) is symmetric.
\end{proof}

\begin{corollary}
The set \(\operatorname{Qis}\) is a multiplicative system in \(\KK(\mathcal{A})\), compatible with the triangulation.
\end{corollary}

\begin{definition}
The \emph{derived category} of \(\mathcal{A}\) is
\[
\DD(\mathcal{A}) := \KK(\mathcal{A})_{\operatorname{Qis}}.
\]
Similarly define
\[
\DD^+(\mathcal{A})=\KK^+(\mathcal{A})_{\operatorname{Qis}},\quad
\DD^-(\mathcal{A})=\KK^-(\mathcal{A})_{\operatorname{Qis}},\quad
\DD^b(\mathcal{A})=\KK^b(\mathcal{A})_{\operatorname{Qis}}.
\]
By Proposition 1.3.3, these are full subcategories of \(\DD(\mathcal{A})\), and
\[
\DD^+(\mathcal{A})\cap \DD^-(\mathcal{A}) = \DD^b(\mathcal{A}).
\]
\end{definition}

\begin{definition}
Let \(\mathcal{A}\) be abelian, and \(\mathcal{A}'\subset\mathcal{A}\) a \emph{thick} full abelian subcategory (closed under extensions in \(\mathcal{A}\)). Define \(\KK_{\mathcal{A}'}(\mathcal{A})\) to be the full subcategory of \(\KK(\mathcal{A})\) consisting of complexes \(X^\bullet\) whose cohomology objects \(H^i(X^\bullet)\) lie in \(\mathcal{A}'\). Since \(\mathcal{A}'\) is thick, \(\KK_{\mathcal{A}'}(\mathcal{A})\) is a triangulated subcategory of \(\KK(\mathcal{A})\).

Define
\[
\DD_{\mathcal{A}'}(\mathcal{A}) := \KK_{\mathcal{A}'}(\mathcal{A})_{\operatorname{Qis}}.
\]
By Proposition 1.3.3, \(\DD_{\mathcal{A}'}(\mathcal{A})\) is the full subcategory of \(\DD(\mathcal{A})\) of complexes with all cohomology in \(\mathcal{A}'\). One defines similarly \(\KK_{\mathcal{A}'}^+(\mathcal{A}), \DD_{\mathcal{A}'}^+(\mathcal{A})\), etc., for bounded-below complexes.
\end{definition}

\begin{definition}
A morphism \(f\colon X^\bullet\to Y^\bullet\) of complexes induces the zero morphism in \(\DD(\mathcal{A})\) if and only if:

\((*)\) There exists \(s\colon Y^\bullet\to Y'^\bullet\in\operatorname{Qis}\) such that \(sf\) is homotopic to zero (equivalently, \(\exists\;t\colon X'^\bullet\to X^\bullet\in\operatorname{Qis}\) such that \(ft\) is homotopic to zero).
\end{definition}

\begin{enumerate}
\item If \(f\) is homotopic to zero, then \((*)\) holds. The converse is false.
Take the complex
\[
X^\bullet:\; 0\to\mathbb{Z}\xrightarrow{2}\mathbb{Z}\to\mathbb{Z}/2\mathbb{Z}\to 0,
\]
and \(f=\operatorname{id}_{X^\bullet}\). Let \(g\colon X^\bullet\to 0\) be the zero map; then \(g\in\operatorname{Qis}\) and \(gf=0\), but \(\operatorname{id}_{X^\bullet}\) is not homotopic to zero.

\item If \(f\) is homotopic to zero, it induces the zero map on cohomology; the converse is false.
Take
\[
X^\bullet,Y^\bullet:\quad 0\to\mathbb{Z}\xrightarrow{2}\mathbb{Z}\to 0.
\]
The identity map induces zero on cohomology, but there is no \(t\in\operatorname{Qis}\) with \(ft\simeq 0\). Any homotopy operator would force \(2k(x)=1\), impossible in \(\mathbb{Z}\).
\end{enumerate}

For morphisms \(f,g\colon X^\bullet\to Y^\bullet\) of complexes, the implications below are all strict:
\[
f=g \;\Longrightarrow\; f\simeq g \;\Longrightarrow\; f=g \text{ in }\DD(\mathcal{A}) \;\Longrightarrow\; H^i(f)=H^i(g)\;\forall i.
\]

\begin{proposition}\label{prop:1.4.3}
Let \(\mathcal{A}\) be identified with the full subcategory of \(\DD(\mathcal{A})\) consisting of complexes concentrated in degree zero. Then the natural embedding \(\mathcal{A}\hookrightarrow\DD(\mathcal{A})\) is fully faithful, i.e., for any \(X,Y\in\operatorname{Ob}\mathcal{A}\),
\[
\operatorname{Hom}_{\mathcal{A}}(X,Y) \cong \operatorname{Hom}_{\DD(\mathcal{A})}(X,Y).
\]
\end{proposition}

We now prove three lemmas, which give an alternative description of \(\DD^+(\mathcal{A})\) when \(\mathcal{A}\) has enough injectives.

\begin{lemma}\label{lem:1.4.4}
Let \(\mathcal{A}\) be abelian, \(f\colon Z^\bullet\to I^\bullet\) a morphism of complexes. Assume:
\begin{enumerate}
\item[(1)] \(Z^\bullet\) is acyclic;
\item[(2)] Each \(I^p\) is injective in \(\mathcal{A}\);
\item[(3)] \(I^\bullet\) is bounded below.
\end{enumerate}
Then \(f\) is homotopic to zero.
\end{lemma}

\begin{lemma}\label{lem:1.4.5}
Let \(\mathcal{A}\) be abelian, \(s\colon I^\bullet\to Y^\bullet\) a morphism of complexes. Assume:
\begin{enumerate}
\item[(1)] \(s\) is a quasi-isomorphism;
\item[(2)] Each \(I^p\) is injective;
\item[(3)] \(I^\bullet\) is bounded below.
\end{enumerate}
Then \(s\) admits a homotopy inverse.
\end{lemma}

\begin{proof}
Form the mapping cone \(Z^\bullet = T(I^\bullet)\oplus Y^\bullet\) of \(s\). Since \(s\) is a quasi-isomorphism, \(Z^\bullet\) is acyclic. The natural map \(v\colon Z^\bullet\to T(I^\bullet)\) satisfies the hypotheses of Lemma 1.4.4, hence is homotopic to zero. Let \((k,t)\colon T(I^\bullet)\oplus Y^\bullet\to I^\bullet\) be the corresponding homotopy. Then
\[
v = (\operatorname{id}_{I^\bullet},0) = (k,t)\circ d_Z + d_I\circ(k,t).
\]
Separating components yields
\[
\operatorname{id}_I = dk + kd + ts,\qquad dt - td = 0.
\]
Thus \(t\colon Y^\bullet\to I^\bullet\) is a morphism of complexes, and \(ts\) is homotopic to \(\operatorname{id}_I\). Hence \(t\) is a homotopy inverse of \(s\).
\end{proof}

\begin{lemma}\label{lem:1.4.6}
Let \(\mathcal{A}\) be an abelian category.
\begin{enumerate}
\item[1)] Let \(P\subseteq\operatorname{Ob}\mathcal{A}\) satisfy:
\begin{enumerate}
\item[(i)] Every object of \(\mathcal{A}\) admits an injection into an element of \(P\).
\end{enumerate}
Then every bounded-below complex \(X^\bullet\in\KK^+(\mathcal{A})\) admits a quasi-isomorphism into a bounded-below complex \(I^\bullet\) of objects of \(P\).

\item[2)] Assume further that \(P\) satisfies:
\begin{enumerate}
\item[(ii)] If \(0\to X\to Y\to Z\to 0\) is exact and \(X\in P\), then \(Y\in P\iff Z\in P\);
\item[(iii)] There exists a positive integer \(n\) such that for any exact sequence
\[
X^0\to X^1\to\cdots\to X^{n-1}\to X^n\to 0,
\]
if \(X^0,\dots,X^{n-1}\in P\) then \(X^n\in P\).
\end{enumerate}
Then every \(X^\bullet\in\KK(\mathcal{A})\) admits a quasi-isomorphism into a complex \(I^\bullet\) of objects of \(P\).

\item[3)] Let \(\mathcal{A}'\) be a thick subcategory of \(\mathcal{A}\), and suppose \(\mathcal{A}'\) has enough \(\mathcal{A}\)-injectives. Then every \(X^\bullet\in\KK_{\mathcal{A}'}^+(\mathcal{A})\) admits a quasi-isomorphism into a bounded-below complex \(I^\bullet\) of \(\mathcal{A}\)-injective objects of \(\mathcal{A}'\).
\end{enumerate}
\end{lemma}

\begin{proof}[Proof of 1)]
We may assume \(X^p=0\) for \(p<0\). Embed \(X^0\hookrightarrow I^0\) with \(I^0\in P\). Having defined \(I^0,I^1,\dots,I^p\), choose \(I^{p+1}\in P\) containing
\[
I^p\big/\operatorname{im}I^{p-1}\;\oplus\;X^{p+1},
\]
and define differentials \(I^p\to I^{p+1}\) and morphisms \(X^{p+1}\to I^{p+1}\) canonically. One verifies \(X^\bullet\to I^\bullet\) is a quasi-isomorphism, and each \(X^p\to I^p\) is injective.
\end{proof}

\begin{proof}[Proof of 2)]
Let \(X^\bullet\) be any complex. For a fixed integer \(i_0\), by (i), the truncated complex
\[
0\to 0\to X^{i_0}\to X^{i_0+1}\to\cdots
\]
admits a quasi-isomorphism into a complex \(I^\bullet\) of objects of \(P\), with each \(X^i\to I^i\) injective. Define \(X_0^\bullet\) by
\[
\cdots\to X^{i_0-2}\to X^{i_0-1}\to I^{i_0}\to I^{i_0+1}\to\cdots.
\]
Then \(X^\bullet\to X_0^\bullet\) is a quasi-isomorphism, \(X_0^i\in P\) for \(i\ge i_0\), and each \(X^i\to X_0^i\) is injective.

Now suppose we have \(X_1^\bullet\) with \(X_1^i\in P\) for \(i\ge i_1\), and let \(i_2<i_1\). We construct a quasi-isomorphism \(X_1^\bullet\to X_2^\bullet\) such that \(X_2^i\in P\) for \(i\ge i_2\) and \(X_1^i=X_2^i\) for \(i\ge i_1+n\) (with \(n\) from condition (iii)).

By the first step, choose a quasi-isomorphism \(X_1^\bullet\to X'^\bullet\) with \(X'^i\in P\) for \(i\ge i_2\) and each map injective. Let \(Y^i=\operatorname{coker}(X_1^i\to X'^i)\). Then \(Y^\bullet\) is acyclic, and \(Y^i\in P\) for \(i\ge i_1\) by (ii). By (iii), the boundaries \(B^i(Y^\bullet)\in P\) for \(i\ge i_1+n\).

Define \(X_2^\bullet\) piecewise:
\[
X_2^i=
\begin{cases}
X'^i & i < i_1+n,\\[2pt]
B^i(X'^\bullet)\oplus \ker(X_1^{i-1}\to X_1^i) & i = i_1+n,\\[2pt]
X_1^i & i > i_1+n.
\end{cases}
\]
From the exact sequence
\[
0\to X_1^i\to B^i(X'^\bullet)\oplus \ker(X_1^{i-1}\to X_1^i)\to B^i(X'^\bullet)\to 0,
\]
condition (ii) implies the middle term lies in \(P\) for \(i\ge i_1+n\). One checks \(X_1^\bullet\to X_2^\bullet\) is a quasi-isomorphism.

Now for any \(X^\bullet\in\KK(\mathcal{A})\), pick a decreasing sequence \(i_0>i_1>i_2>\cdots\to-\infty\). Construct \(X_0^\bullet,X_1^\bullet,X_2^\bullet,\dots\) as above, giving quasi-isomorphisms
\[
X^\bullet\to X_0^\bullet\to X_1^\bullet\to X_2^\bullet\to\cdots.
\]
For each fixed degree \(i\), the sequence \(X^i\to X_0^i\to X_1^i\to\cdots\) is eventually constant and lies in \(P\). The direct limit \(I^\bullet=\varinjlim X_r^\bullet\) is the desired complex of objects of \(P\).
\end{proof}

\begin{proof}[Proof of 3)]
Assume \(X^i=0\) for \(i<0\). Embed \(H^0(X^\bullet)\hookrightarrow I^0\), where \(I^0\) is an \(\mathcal{A}\)-injective object of \(\mathcal{A}'\); this is possible since \(H^0(X^\bullet)\in\mathcal{A}'\) and \(\mathcal{A}'\) has enough \(\mathcal{A}\)-injectives. Extend to a morphism \(f^0\colon X^0\to I^0\).

Inductively having defined \(I^0,\dots,I^p\) and \(f^i\colon X^i\to I^i\) for \(0\le i\le p\), choose an \(\mathcal{A}\)-injective \(I^{p+1}\in\mathcal{A}'\) containing
\[
I^p\big/\operatorname{im}I^{p-1}\;\oplus\;Z^{p+1}(X^\bullet).
\]
Since \(\mathcal{A}'\) is thick, this object belongs to \(\mathcal{A}'\). Extend the natural map \(Z^{p+1}(X^\bullet)\to I^{p+1}\) to \(f^{p+1}\colon X^{p+1}\to I^{p+1}\). The resulting morphism \(f\colon X^\bullet\to I^\bullet\) is a quasi-isomorphism.
\end{proof}

\begin{proposition}\label{prop:1.4.7}
Let \(\mathcal{A}\) be abelian, and let \(\mathcal{I}\) be the additive subcategory of injective objects of \(\mathcal{A}\). The natural functor
\[
\alpha\colon \KK^+(\mathcal{I})\to \DD^+(\mathcal{A})
\]
is fully faithful. If furthermore \(\mathcal{A}\) has enough injectives, then \(\alpha\) is an equivalence of categories.
\end{proposition}

\begin{proof}
By Proposition 1.4.2, \(\operatorname{Qis}\cap\KK^+(\mathcal{I})\) is a multiplicative system in \(\KK^+(\mathcal{I})\). Lemma 1.4.5 verifies condition (ii) of Proposition 1.3.3 for the inclusion \(\KK^+(\mathcal{I})\subseteq\KK^+(\mathcal{A})\) and \(\operatorname{Qis}\). Thus
\[
\DD^+(\mathcal{I})\to\DD^+(\mathcal{A})
\]
is fully faithful. By Lemma 1.4.5, every quasi-isomorphism in \(\KK^+(\mathcal{I})\) is an isomorphism, so \(\KK^+(\mathcal{I}) = \DD^+(\mathcal{I})\).

If \(\mathcal{A}\) has enough injectives, apply Lemma 1.4.6(1) with \(P=\operatorname{Ob}\mathcal{I}\): every object of \(\DD^+(\mathcal{A})\) is isomorphic to an object of \(\KK^+(\mathcal{I})\). Hence \(\alpha\) is an equivalence.
\end{proof}

\begin{proposition}\label{prop:1.4.8}
Let \(\mathcal{A}\) be abelian, \(\mathcal{A}'\subseteq\mathcal{A}\) a thick abelian subcategory. Assume \(\mathcal{A}'\) has enough \(\mathcal{A}\)-injectives. Then the natural functor
\[
\DD^+(\mathcal{A}') \to \DD_{\mathcal{A}'}^+(\mathcal{A})
\]
is an equivalence of categories.
\end{proposition}

\begin{proof}
Apply Proposition 1.3.3 to the inclusion \(\KK^+(\mathcal{A}')\hookrightarrow\KK^+(\mathcal{A})\). The set \(\operatorname{Qis}\) is a multiplicative system in both categories. Any quasi-isomorphism \(X^\bullet\to Y^\bullet\) with \(X^\bullet\in\KK^+(\mathcal{A}')\) has cohomology in \(\mathcal{A}'\); by Lemma 1.4.6(3), \(Y^\bullet\) admits a quasi-isomorphism into an injective complex \(I^\bullet\in\KK^+(\mathcal{A}')\). Thus condition (ii) of Proposition 1.3.3 holds, so \(\DD^+(\mathcal{A}')\to\DD^+(\mathcal{A})\) is fully faithful, with image exactly \(\DD_{\mathcal{A}'}^+(\mathcal{A})\).
\end{proof}

\textbf{Exercise.} Formulate the analogous statements for projective objects of \(\mathcal{A}\) and the bounded-above derived category \(\DD^-(\mathcal{A})\).

\section{Derived functors}

We treat only \emph{right derived covariant functors}; the reader may make the obvious modifications for left derived covariant functors, and right/left derived contravariant functors.

Let \(\mathcal{A},\mathcal{B}\) be abelian categories, and let \(F\colon\KK(\mathcal{A})\to\KK(\mathcal{B})\) be a triangulated functor (cf.\ \S1). This holds in particular for any additive functor \(F\colon\mathcal{A}\to\mathcal{B}\), which extends canonically to \(\KK(\mathcal{A})\).

In general, \(F\) need not send quasi-isomorphisms to quasi-isomorphisms; if it does, then \(F\) localizes to induce a functor \(\DD(\mathcal{A})\to\DD(\mathcal{B})\). Exact functors are examples of this behavior.

We are led to ask whether there exists a functor \(\DD(\mathcal{A})\to\DD(\mathcal{B})\) approximating \(F\); this motivates the general definition of derived functors below.

\begin{definition}
Let \(\mathcal{A}\) be abelian, and let \(\KK^*(\mathcal{A})\subseteq\KK(\mathcal{A})\) be a triangulated subcategory. By Proposition 1.4.2, \(\KK^*(\mathcal{A})\cap\operatorname{Qis}\) is a multiplicative system in \(\KK^*(\mathcal{A})\).

We say \(\KK^*(\mathcal{A})\) is a \emph{localizing subcategory} of \(\KK(\mathcal{A})\) if the natural functor
\[
\KK^*(\mathcal{A})_{\KK^*(\mathcal{A})\cap\operatorname{Qis}}
\;\to\;
\KK(\mathcal{A})_{\operatorname{Qis}} = \DD(\mathcal{A})
\]
is fully faithful. In this case we write \(\DD^*(\mathcal{A})\) for the domain category above.
\end{definition}

\textbf{Examples.}
\begin{enumerate}
\item The intersection of any family of localizing subcategories is localizing.
\item \(\KK^+(\mathcal{A}),\KK^-(\mathcal{A}),\KK^b(\mathcal{A})\) are localizing subcategories of \(\KK(\mathcal{A})\) (cf.\ \S4).
\item If \(\mathcal{A}'\subseteq\mathcal{A}\) is a thick subcategory, then \(\KK_{\mathcal{A}'}(\mathcal{A}),\KK_{\mathcal{A}'}^+(\mathcal{A}),\KK_{\mathcal{A}'}^-(\mathcal{A}),\KK_{\mathcal{A}'}^b(\mathcal{A})\) are localizing subcategories of \(\KK(\mathcal{A})\) (cf.\ \S4).
\item The complexes of finite injective dimension form a localizing subcategory \(\KK^+(\mathcal{A})_{\mathrm{fid}}\subseteq\KK(\mathcal{A})\) (cf.\ Corollary 1.7.7).
\end{enumerate}

\begin{definition}
Let \(\mathcal{A},\mathcal{B}\) be abelian categories, \(\KK^*(\mathcal{A})\subseteq\KK(\mathcal{A})\) a localizing subcategory, and
\[
F\colon\KK^*(\mathcal{A})\to\KK(\mathcal{B})
\]
a triangulated functor.

Denote by \(Q\) the localization functor \(\KK^*(\mathcal{A})\to\DD^*(\mathcal{A})\) (resp.\ \(\KK(\mathcal{B})\to\DD(\mathcal{B})\)). The \emph{right derived functor} of \(F\) is a pair
\[
\mathbf{R}^*F\colon\DD^*(\mathcal{A})\to\DD(\mathcal{B})
\]
together with a functor morphism
\[
\xi\colon Q\circ F \;\to\; \mathbf{R}^*F\circ Q
\]
of functors \(\KK^*(\mathcal{A})\to\DD(\mathcal{B})\), satisfying the following universal property:

For any triangulated functor \(G\colon\DD^*(\mathcal{A})\to\DD(\mathcal{B})\) and any functor morphism
\[
\zeta\colon Q\circ F \;\to\; G\circ Q,
\]
there exists a unique functor morphism
\[
\eta\colon \mathbf{R}^*F \to G
\]
such that
\[
\zeta = (\eta\circ Q)\circ \xi.
\]
\end{definition}

\textbf{Remarks.}
\begin{enumerate}
\item If \(\mathbf{R}^*F\) exists, it is unique up to unique isomorphism of functors.
\item For standard localizing subcategories \(\KK^+(\mathcal{A}),\KK^-(\mathcal{A}),\KK_{\mathcal{A}'}(\mathcal{A})\), we write
\[
\mathbf{R}^+F,\quad \mathbf{R}^-F,\quad \mathbf{R}_{\mathcal{A}'}F
\]
respectively; when no confusion arises, we write simply \(\mathbf{R}F\).

\item We define \(R^iF = H^i(\mathbf{R}F)\). If \(F\) is left-exact on \(\mathcal{A}\) and \(\mathcal{A}\) has enough injectives, these coincide with the classical right derived functors of \(F\).
\item Let \(\varphi\colon F\to G\) be a morphism of functors \(\KK^*(\mathcal{A})\to\KK(\mathcal{B})\). If \(\mathbf{R}F,\mathbf{R}G\) exist, there is a unique induced morphism
\[
\mathbf{R}\varphi\colon \mathbf{R}F\to\mathbf{R}G
\]
compatible with the structure morphisms \(\xi\). This follows directly from the universal property.
\item Let \(\KK^{**}(\mathcal{A})\subseteq\KK^*(\mathcal{A})\) be localizing subcategories, \(F\colon\KK^*(\mathcal{A})\to\KK(\mathcal{B})\) a triangulated functor. If both \(\mathbf{R}^*F\) and \(\mathbf{R}^{**}(F|_{\KK^{**}(\mathcal{A})})\) exist, there is a natural functor morphism
\[
\mathbf{R}^{**}\big(F|_{\KK^{**}(\mathcal{A})}\big)
\;\to\;
\mathbf{R}^*F\big|_{\DD^{**}(\mathcal{A})}.
\]
We do not know if this is an isomorphism in general, but it holds in all our intended applications (cf.\ Corollary 1.5.3).
\end{enumerate}

\begin{theorem}[Existence of Derived Functors]\label{thm:1.5.1}
Let \(\mathcal{A},\mathcal{B},\KK^*(\mathcal{A}),F\) be as above. Suppose there exists a triangulated subcategory \(\LL\subseteq\KK^*(\mathcal{A})\) such that:
\begin{enumerate}
\item[(1)] Every object of \(\KK^*(\mathcal{A})\) admits a quasi-isomorphism into some object of \(\LL\);
\item[(2)] For every acyclic complex \(I^\bullet\in\operatorname{Ob}\LL\) (i.e.\ \(H^i(I^\bullet)=0\) for all \(i\)), the complex \(F(I^\bullet)\) is also acyclic.
\end{enumerate}
Then the right derived functor \((\mathbf{R}^*F,\xi)\) exists. Furthermore, for any \(I^\bullet\in\operatorname{Ob}\LL\), the structure morphism
\[
\xi(I^\bullet)\colon Q\circ F(I^\bullet) \to \mathbf{R}^*F\circ Q(I^\bullet)
\]
is an isomorphism in \(\DD(\mathcal{B})\).
\end{theorem}

\begin{proof}
First, the restriction \(F|_\LL\) sends quasi-isomorphisms to quasi-isomorphisms. Indeed, let \(s\colon I_1^\bullet\to I_2^\bullet\) be a quasi-isomorphism in \(\LL\); complete \(s\) to a triangle in \(\KK^*(\mathcal{A})\). The third vertex \(J^\bullet\) of the triangle is acyclic, so by hypothesis (2), \(F(J^\bullet)\) is acyclic. Hence \(F(s)\) is a quasi-isomorphism.

Thus \(F\) descends to a functor
\[
\overline{F}\colon\LL_{\operatorname{Qis}} \to \DD(\mathcal{B})
\]
satisfying \(\overline{F}\circ Q = Q\circ F\) on \(\LL\) (with \(Q\) denoting localization functors as usual).

Second, the pair \((\LL,\KK^*(\mathcal{A}),\operatorname{Qis})\) satisfies condition (ii) of Proposition 1.3.3, so the natural functor
\[
T\colon\LL_{\operatorname{Qis}} \to \DD^*(\mathcal{A})
\]
is an equivalence of categories. Let \(U\colon\DD^*(\mathcal{A})\to\LL_{\operatorname{Qis}}\) be a quasi-inverse of \(T\), with functorial isomorphisms
\[
\alpha\colon \operatorname{id}_{\LL_{\operatorname{Qis}}} \to U\circ T,\qquad
\beta\colon \operatorname{id}_{\DD^*(\mathcal{A})} \to T\circ U.
\]

Define
\[
\mathbf{R}^*F := \overline{F}\circ U.
\]

We now define the functor morphism
\[
\xi\colon Q\circ F \to \mathbf{R}^*F\circ Q = \overline{F}\circ U\circ Q
\]
as follows. Given \(X^\bullet\in\operatorname{Ob}\KK^*(\mathcal{A})\), choose \(I^\bullet\in\operatorname{Ob}\LL\) such that \(Q(I^\bullet)\cong U\circ Q(X^\bullet)\) in \(\DD^*(\mathcal{A})\).

We have the canonical isomorphism
\[
\beta\big(Q(X^\bullet)\big)\colon Q(X^\bullet) \stackrel{\sim}{\longrightarrow}
T\circ U\big(Q(X^\bullet)\big) = T\big(Q(I^\bullet)\big).
\]
This isomorphism may be represented by a diagram in \(\KK^*(\mathcal{A})\)
\[
X^\bullet \xleftarrow{s} Y^\bullet \xrightarrow{t} I^\bullet
\]
with \(s,t\in\operatorname{Qis}\) and \(Y^\bullet\in\operatorname{Ob}\KK^*(\mathcal{A})\). By hypothesis (1), we may assume \(Y^\bullet\in\operatorname{Ob}\LL\). Applying the triangulated functor \(F\), we obtain a diagram in \(\KK(\mathcal{B})\):
\[
F(X^\bullet) \xleftarrow{F(s)} F(Y^\bullet) \xrightarrow{F(t)} F(I^\bullet).
\]
As observed previously, \(F(s),F(t)\) are quasi-isomorphisms.

This induces a morphism in \(\DD(\mathcal{B})\):
\[
\xi(X^\bullet)\colon Q\circ F(X^\bullet)
\to Q\circ F(I^\bullet)
= \overline{F}\circ Q(I^\bullet)
= \overline{F}\circ U\circ Q(X^\bullet)
= \mathbf{R}^*F\circ Q(X^\bullet).
\]

One verifies that \(\xi(X^\bullet)\) is independent of the choice of diagram above, that \(\xi\) defines a functor morphism \(Q\circ F\to \mathbf{R}^*F\circ Q\), and that the pair \((\mathbf{R}^*F,\xi)\) satisfies the universal property of the right derived functor.

If \(X^\bullet\in\operatorname{Ob}\LL\), then \(F(t)\) is also a quasi-isomorphism in the construction above, so \(\xi(X^\bullet)\) is an isomorphism in \(\DD(\mathcal{B})\), as required.
\end{proof}

\begin{proposition}\label{prop:1.5.2}
Let \(\mathcal{A},\mathcal{B},\KK^*(\mathcal{A}),F\) be as above, and let \(\KK^{**}(\mathcal{A})\subseteq\KK^*(\mathcal{A})\) be another localizing subcategory of \(\KK(\mathcal{A})\). Suppose there exists a triangulated subcategory \(\LL\subseteq\KK^*(\mathcal{A})\) satisfying hypotheses (1),(2) of Theorem 1.5.1, and such that \(\LL\cap\KK^{**}(\mathcal{A})\) satisfies hypothesis (1) for \(\KK^{**}(\mathcal{A})\).
If \(\mathbf{R}^*F\) and \(\mathbf{R}^{**}\big(F|_{\KK^{**}(\mathcal{A})}\big)\) both exist, then the natural morphism
\[
\mathbf{R}^{**}\big(F|_{\KK^{**}(\mathcal{A})}\big)
\;\to\;
\mathbf{R}^*F\big|_{\DD^{**}(\mathcal{A})}
\]
is an isomorphism.
\end{proposition}

\begin{proof}
Existence follows immediately from Theorem 1.5.1. For the isomorphism, any object of \(\DD^{**}(\mathcal{A})\) is isomorphic to \(Q(I^\bullet)\) for some \(I^\bullet\in\LL\cap\KK^{**}(\mathcal{A})\). The claim then follows from the final assertion of Theorem 1.5.1.
\end{proof}

\begin{corollary}\label{cor:1.5.3}
\emph{\(\alpha\))} Let \(\mathcal{A},\mathcal{B}\) be abelian categories, \(F\colon\KK^+(\mathcal{A})\to\KK(\mathcal{B})\) a triangulated functor (e.g., induced by an additive functor \(F_0\colon\mathcal{A}\to\mathcal{B}\)). If \(\mathcal{A}\) has enough injectives, then \(\mathbf{R}^+F\) exists.

\emph{\(\beta\))} Let \(\mathcal{A},\mathcal{B}\) be abelian categories, \(F\colon\mathcal{A}\to\mathcal{B}\) additive. Suppose there exists \(P\subseteq\operatorname{Ob}\mathcal{A}\) satisfying (i),(ii) of Lemma 1.4.6, and:
\begin{enumerate}
\item[(iv)] \(F\) sends short exact sequences in \(\mathcal{A}\) to short exact sequences in \(\mathcal{B}\).
\end{enumerate}
Then \(\mathbf{R}^+F\) exists (we also denote by \(F\) its canonical extension \(\KK^+(\mathcal{A})\to\KK^+(\mathcal{B})\)).

\emph{\(\gamma\))} Let \(\mathcal{A},\mathcal{B}\) be abelian categories, \(F\colon\mathcal{A}\to\mathcal{B}\) additive. Assume:
\begin{enumerate}
\item[(a)] The hypotheses of \(\beta\) hold;
\item[(b)] \(F\) has finite cohomological dimension: there exists \(n>0\) such that \(R^iF(Y)=0\) for all \(Y\in\operatorname{Ob}\mathcal{A}\) and all \(i>n\).
\end{enumerate}
Then \(\mathbf{R}F\) exists on all of \(\DD(\mathcal{A})\), and its restriction to \(\DD^+(\mathcal{A})\) coincides with \(\mathbf{R}^+F\).
\end{corollary}

\textbf{Remark.} Case \(\alpha\) is a special case of \(\beta\): if \(\mathcal{A}\) has enough injectives, the class of injective objects satisfies (i),(ii),(iv) for any additive \(F\).

\begin{proof}
\emph{\(\alpha\))} Let \(\LL\subseteq\KK^+(\mathcal{A})\) be the triangulated subcategory of bounded-below injective complexes. By Lemma 1.4.6(1), every object of \(\KK^+(\mathcal{A})\) admits a quasi-isomorphism into an object of \(\LL\). By Lemma 1.4.5, every quasi-isomorphism in \(\LL\) is an isomorphism, so acyclic objects of \(\LL\) are zero in \(\KK(\mathcal{A})\). Thus condition (2) of Theorem 1.5.1 holds for any triangulated \(F\), so \(\mathbf{R}^+F\) exists.

\emph{\(\beta\))} Take \(\LL\subseteq\KK^+(\mathcal{A})\) to be complexes of objects of \(P\); \(\LL\) is triangulated by stability of \(P\) under extensions. Condition (1) holds by Lemma 1.4.6(1). For (2), let \(Z^\bullet\in\LL\) be acyclic; using (ii) repeatedly, all kernels \(\ker d_Z^p\) lie in \(P\). By (iv), \(F\) preserves exactness, so \(F(Z^\bullet)\) is acyclic. Theorem 1.5.1 applies, hence \(\mathbf{R}^+F\) exists.

\emph{\(\gamma\))} Let \(P'\) be the class of \(F\)-acyclic objects of \(\mathcal{A}\) (\(R^iF(X)=0,\forall i>0\)). Then \(P'\) satisfies (i),(ii),(iii) of Lemma 1.4.6. Let \(\LL\subseteq\KK(\mathcal{A})\) be complexes of objects of \(P'\). By Lemma 1.4.6 and the argument for \(\beta\), the hypotheses of Theorem 1.5.1 are satisfied, so \(\mathbf{R}F\) exists.

By Proposition 1.5.2, the restriction of \(\mathbf{R}F\) to \(\DD^+(\mathcal{A})\) is isomorphic to \(\mathbf{R}^+F\).
\end{proof}

\begin{proposition}\label{prop:1.5.4}
Let \(\mathcal{A},\mathcal{B},\mathcal{C}\) be abelian categories, \(\KK^*(\mathcal{A})\subseteq\KK(\mathcal{A})\), \(\KK^\dagger(\mathcal{B})\subseteq\KK(\mathcal{B})\) localizing subcategories, and let
\[
F\colon\KK^*(\mathcal{A})\to\KK(\mathcal{B}),\qquad
G\colon\KK^\dagger(\mathcal{B})\to\KK(\mathcal{C})
\]
be triangulated functors.

\emph{a)} Assume \(F(\KK^*(\mathcal{A}))\subseteq\KK^\dagger(\mathcal{B})\), that \(\mathbf{R}^*F,\mathbf{R}^\dagger G,\mathbf{R}^*(G\circ F)\) exist, and \(\mathbf{R}^*F(\DD^*(\mathcal{A}))\subseteq\DD^\dagger(\mathcal{B})\). Then there exists a unique functor morphism
\[
\zeta\colon \mathbf{R}^*(G\circ F) \to \mathbf{R}^\dagger G\circ \mathbf{R}^*F
\]
making the diagram of structure maps \(\xi\) commutative.

\emph{b)} Suppose there exist triangulated subcategories \(\LL\subseteq\KK^*(\mathcal{A})\), \(\MM\subseteq\KK^\dagger(\mathcal{B})\) satisfying Theorem 1.5.1's hypotheses for \(F,G\) respectively, and \(F(\LL)\subseteq\MM\). Then the hypotheses of (a) hold, and the canonical morphism \(\zeta\) is an isomorphism.
\end{proposition}

\begin{proof}
Straightforward from universal properties.
\end{proof}

\textbf{Remarks.}
\begin{enumerate}
\item For three composable functors \(F,G,H\), the morphisms \(\zeta\) yield a commutative diagram of derived functor compositions whenever all derived functors exist.
\item This proposition illustrates the power of derived categories: classical composite functor spectral sequences are replaced by direct functor composition in the derived category. The spectral sequence may be recovered by taking cohomology of the derived functor morphism.
\end{enumerate}

\begin{proposition}\label{prop:1.5.6}
Let \(\mathcal{A}'\subseteq\mathcal{A}\) be a thick abelian subcategory, \(\mathcal{B}\) another abelian category, \(F\colon\KK^+(\mathcal{A})\to\KK^+(\mathcal{B})\) a triangulated functor. Suppose \(\mathbf{R}^+F\) and \(\mathbf{R}^+\big(F|_{\mathcal{A}'}\big)\) both exist. There is a natural functor morphism
\[
\zeta\colon \mathbf{R}^+\big(F|_{\mathcal{A}'}\big) \to \varphi\circ \mathbf{R}^+F,
\]
where \(\varphi\colon\DD^+(\mathcal{A}')\to\DD^+(\mathcal{A})\) is the natural embedding. If additionally \(\mathcal{A}'\) has enough \(\mathcal{A}\)-injectives and \(\mathcal{A}\) has enough injectives, then \(\zeta\) is an isomorphism.
\end{proposition}

\begin{proof}
Existence of \(\zeta\) follows from the definition of derived functors. If \(\mathcal{A}'\) has enough \(\mathcal{A}\)-injectives, we may compute both derived functors via injective resolutions in \(\mathcal{A}'\); by Proposition 1.4.8, \(\varphi\) is fully faithful, hence \(\zeta\) is an isomorphism.
\end{proof}

\section{Examples. \(\operatorname{Ext}\) and \(\mathbf{R}\mathcal{H}\textit{om}\)}

\begin{definition}
Let \(\mathcal{A}\) be an abelian category, and let \(X^\bullet,Y^\bullet\in\DD(\mathcal{A})\).
Define the \emph{hyperext groups} by
\[
\operatorname{Ext}^i\big(X^\bullet,Y^\bullet\big)
= \operatorname{Hom}_{\DD(\mathcal{A})}\big(X^\bullet, T^i(Y^\bullet)\big).
\]
\end{definition}

\textbf{Remarks.}
\begin{enumerate}
\item If \(X^\bullet,Y^\bullet\in\DD^+(\mathcal{A})\), then \(\operatorname{Ext}^i\) computed in \(\DD^+(\mathcal{A})\)
agrees with the above, since \(\DD^+(\mathcal{A})\) is a full subcategory of \(\DD(\mathcal{A})\).

\item This definition gives \(\operatorname{Ext}^i(X,Y)\) for any \(X,Y\in\mathcal{A}\).
We shall see that if \(\mathcal{A}\) has enough injectives, this coincides with the classical
definition of Ext groups.
\end{enumerate}

\begin{proposition}\label{prop:1.6.1}
Let
\[
0 \to X^\bullet \to Y^\bullet \to Z^\bullet \to 0
\]
be a short exact sequence of complexes over \(\mathcal{A}\), and let \(V^\bullet\) be another complex.
Then we have long exact sequences
\[
\cdots \to \operatorname{Ext}^i(V^\bullet,X^\bullet)
\to \operatorname{Ext}^i(V^\bullet,Y^\bullet)
\to \operatorname{Ext}^i(V^\bullet,Z^\bullet)
\to \operatorname{Ext}^{i+1}(V^\bullet,X^\bullet)
\to \cdots
\]
and
\[
\cdots \to \operatorname{Ext}^i(Z^\bullet,V^\bullet)
\to \operatorname{Ext}^i(Y^\bullet,V^\bullet)
\to \operatorname{Ext}^i(X^\bullet,V^\bullet)
\to \operatorname{Ext}^{i+1}(Z^\bullet,V^\bullet)
\to \cdots
\]
\end{proposition}

\begin{proof}
Complete \(X^\bullet\to Y^\bullet\) to a triangle in \(\KK(\mathcal{A})\), with third vertex \(W^\bullet\):
\[
X^\bullet \to Y^\bullet \to W^\bullet \to T(X^\bullet).
\]
There exists a morphism \(g\colon W^\bullet\to Z^\bullet\). Using the long exact cohomology sequence
and the five-lemma, one verifies \(g\) is a quasi-isomorphism, hence an isomorphism in \(\DD(\mathcal{A})\).
Replacing \(Z^\bullet\) by \(W^\bullet\), the result follows from properties of triangles.
\end{proof}

\textbf{Remark.}
Any short exact sequence of complexes
\[
0 \to X^\bullet \to Y^\bullet \to Z^\bullet \to 0
\]
determines a morphism \(Z^\bullet\to T(X^\bullet)\) in \(\DD(\mathcal{A})\), making \(X^\bullet,Y^\bullet,Z^\bullet\)
into a triangle.

We now define a complex whose cohomology yields the \(\operatorname{Ext}\) groups.

\begin{definition}
For complexes \(X^\bullet,Y^\bullet\) over \(\mathcal{A}\), define the complex
\(\mathcal{H}\textit{om}^\bullet(X^\bullet,Y^\bullet)\) by
\[
\mathcal{H}\textit{om}^n\big(X^\bullet,Y^\bullet\big)
= \prod_{p\in\mathbb{Z}} \operatorname{Hom}_{\mathcal{A}}\big(X^p,Y^{p+n}\big),
\]
with differential
\[
d^n = \prod_{p\in\mathbb{Z}} \big(d_X^{p-1} + (-1)^{n+1}d_Y^{p+n}\big).
\]

The \(n\)-cycles of \(\mathcal{H}\textit{om}^\bullet(X^\bullet,Y^\bullet)\) correspond bijectively
to morphisms \(X^\bullet\to T^n(Y^\bullet)\) of complexes.
The \(n\)-boundaries correspond to morphisms homotopic to zero.
Thus
\[
H^n\big(\mathcal{H}\textit{om}^\bullet(X^\bullet,Y^\bullet)\big)
\cong \operatorname{Hom}_{\KK(\mathcal{A})}\big(X^\bullet, T^n(Y^\bullet)\big).
\]

Clearly \(\mathcal{H}\textit{om}^\bullet\) is a bifunctor
\[
\mathcal{H}\textit{om}^\bullet\colon\KK(\mathcal{A})^\circ \times \KK(\mathcal{A}) \to \KK(\mathbf{Ab}).
\]
\end{definition}

If \(\mathcal{A}\) has enough injectives, we may form its derived functors.

\begin{lemma}\label{lem:1.6.2}
Let \(X^\bullet\in\KK(\mathcal{A})\), \(Y^\bullet\in\KK^+(\mathcal{A})\) be a bounded-below complex of injectives.
Assume either
\begin{enumerate}
\item[(a)] \(Y^\bullet\) is acyclic, or
\item[(b)] \(X^\bullet\) is acyclic.
\end{enumerate}
Then \(\mathcal{H}\textit{om}^\bullet(X^\bullet,Y^\bullet)\) is acyclic.
\end{lemma}

\begin{proof}
It suffices to show any morphism \(X^\bullet\to Y^\bullet\) is homotopic to zero.
Case (a): An acyclic bounded-below injective complex is split exact; one constructs a homotopy directly.
Case (b): Follows from Lemma 1.4.4.
\end{proof}

Suppose \(\mathcal{A}\) has enough injectives, and let \(\LL\subseteq\KK^+(\mathcal{A})\) be the triangulated
subcategory of bounded-below injective complexes.
By Lemma 1.6.2(a), for fixed \(X^\bullet\in\KK(\mathcal{A})\), the functor
\[
\mathcal{H}\textit{om}^\bullet(X^\bullet,\,\cdot\,)\colon\KK^+(\mathcal{A})\to\KK(\mathbf{Ab})
\]
satisfies the hypotheses of Theorem 1.5.1, hence admits a right derived functor.
This is functorial in \(X^\bullet\), giving a biderived functor
\[
\mathbf{R}_\mathrm{II}\mathcal{H}\textit{om}^\bullet
\colon\KK(\mathcal{A})^\circ \times \DD^+(\mathcal{A}) \to \DD(\mathbf{Ab}).
\]

By Lemma 1.6.2(b), this functor sends acyclic complexes to acyclic complexes in the first variable,
hence descends to
\[
\mathbf{R}\mathcal{H}\textit{om}^\bullet
\colon\DD(\mathcal{A})^\circ \times \DD^+(\mathcal{A}) \to \DD(\mathbf{Ab}).
\]

Dually, if \(\mathcal{A}\) has enough projectives, we obtain
\[
\mathbf{R}_\mathrm{I}\mathcal{H}\textit{om}^\bullet
\colon\DD^-(\mathcal{A})^\circ \times \DD(\mathcal{A}) \to \DD(\mathbf{Ab}).
\]

If \(\mathcal{A}\) has both enough injectives and projectives, the two derived functors
are canonically isomorphic on \(\DD^-(\mathcal{A})^\circ\times\DD^+(\mathcal{A})\).
We therefore write simply \(\mathbf{R}\mathcal{H}\textit{om}^\bullet\).

\begin{lemma}\label{lem:1.6.3}
Let \(\mathcal{A},\mathcal{B},\mathcal{C}\) be abelian categories,
\[
T\colon\KK^*(\mathcal{A})\times\KK^\dagger(\mathcal{B})\to\KK(\mathcal{C})
\]
a bitriangulated functor.
If the iterated derived functors \(\mathbf{R}_\mathrm{I}\mathbf{R}_\mathrm{II}T\) and \(\mathbf{R}_\mathrm{II}\mathbf{R}_\mathrm{I}T\) both exist,
they are canonically isomorphic.
\end{lemma}

\begin{proof}
Immediate from the universal property of derived functors.
\end{proof}

\begin{theorem}[Yoneda]\label{thm:1.6.4}
Let \(\mathcal{A}\) have enough injectives.
For all \(X^\bullet\in\DD(\mathcal{A})\), \(Y^\bullet\in\DD^+(\mathcal{A})\),
\[
H^i\big(\mathbf{R}\mathcal{H}\textit{om}^\bullet(X^\bullet,Y^\bullet)\big)
= \operatorname{Ext}^i\big(X^\bullet,Y^\bullet\big).
\]
\end{theorem}

\begin{corollary}\label{cor:1.6.5}
If \(\mathcal{A}\) has enough injectives, the derived-category definition of \(\operatorname{Ext}^i(X,Y)\)
coincides with the classical Ext groups for all \(X,Y\in\mathcal{A}\).
\end{corollary}

\begin{proof}[Proof of Corollary 1.6.5]
Take an injective resolution \(I^\bullet\) of \(Y\). Then
\[
\operatorname{Ext}^i(X,Y) = H^i\big(\mathcal{H}\textit{om}^\bullet(X,I^\bullet)\big),
\]
which is the standard definition.
\end{proof}

\begin{proof}[Proof of Theorem 1.6.4]
Choose a quasi-isomorphism \(s\colon Y^\bullet\to I^\bullet\) with \(I^\bullet\) a bounded-below injective complex. Then
\[
\operatorname{Ext}^i(X^\bullet,Y^\bullet)
= \operatorname{Ext}^i(X^\bullet,I^\bullet)
= \operatorname{Hom}_{\DD(\mathcal{A})}\big(X^\bullet,T^i(I^\bullet)\big).
\]
By Lemma 1.4.5, any morphism \(X^\bullet\to T^i(I^\bullet)\) in \(\DD(\mathcal{A})\)
is represented by an actual morphism of complexes. Therefore
\[
\begin{aligned}
\operatorname{Hom}_{\DD(\mathcal{A})}\big(X^\bullet,T^i(I^\bullet)\big)
&= \operatorname{Hom}_{\KK(\mathcal{A})}\big(X^\bullet,T^i(I^\bullet)\big) \\
&= H^i\big(\mathcal{H}\textit{om}^\bullet(X^\bullet,I^\bullet)\big) \\
&= H^i\big(\mathbf{R}\mathcal{H}\textit{om}^\bullet(X^\bullet,Y^\bullet)\big).
\end{aligned}
\]
\end{proof}

\section{Way-out functors}

\begin{definition}
Let \(\mathcal{A},\mathcal{B}\) be abelian categories, and let \(F\colon\DD(\mathcal{A})\to\DD(\mathcal{B})\) be a triangulated functor.
We say \(F\) is \emph{way-out right} if for every \(n_1\in\mathbb{Z}\), there exists \(n_2\in\mathbb{Z}\) such that
for any complex \(X^\bullet\in\DD(\mathcal{A})\) satisfying \(H^i(X^\bullet)=0\) for all \(i<n_2\),
we have \(H^i(F(X^\bullet))=0\) for all \(i<n_1\).

One defines \emph{way-out left} and \emph{way-out in both directions} similarly.
For a contravariant triangulated functor \(F\colon\DD(\mathcal{A})\to\DD(\mathcal{B})\),
the definition of way-out right is the same, except reversing the inequality to \(i>n_2\).
\end{definition}

\textbf{Examples.}
\begin{enumerate}
\item If \(F_0\colon\mathcal{A}\to\mathcal{B}\) is additive and satisfies the hypotheses of Corollary 1.5.3(\(\alpha\)) or (\(\beta\)),
then \(\mathbf{R}^+F\) is way-out right.
If \(F_0\) satisfies (\(\gamma\)), then \(\mathbf{R}F\) is way-out in both directions.

\item For unbounded \(X^\bullet\in\DD(\mathcal{A})\), the functor
\(\mathbf{R}\mathcal{H}\textit{om}(X^\bullet,\,\cdot\,)\) on \(\DD^+(\mathcal{A})\)
is generally not way-out.
\end{enumerate}

\begin{proposition}[Lemma on Way-out Functors]\label{prop:1.7.1}
Let \(\mathcal{A},\mathcal{B}\) be abelian categories, \(\mathcal{A}'\subseteq\mathcal{A}\) a thick subcategory.
Let \(F,G\colon\DD_{\mathcal{A}'}^+(\mathcal{A})\to\DD(\mathcal{B})\) be triangulated functors,
and \(\eta\colon F\to G\) a morphism of functors.
\begin{enumerate}
\item[(i)] If \(\eta(X)\) is an isomorphism for all \(X\in\operatorname{Ob}\mathcal{A}'\),
then \(\eta(X^\bullet)\) is an isomorphism for all \(X^\bullet\in\DD_{\mathcal{A}'}^b(\mathcal{A})\).

\item[(ii)] If \(\eta(X)\) is an isomorphism for all \(X\in\operatorname{Ob}\mathcal{A}'\),
and \(F,G\) are way-out right,
then \(\eta(X^\bullet)\) is an isomorphism for all \(X^\bullet\in\DD_{\mathcal{A}'}^+(\mathcal{A})\).

\item[(iii)] If \(\eta(X)\) is an isomorphism for all \(X\in\operatorname{Ob}\mathcal{A}'\),
and \(F,G\) are way-out in both directions,
then \(\eta(X^\bullet)\) is an isomorphism for all \(X^\bullet\in\DD_{\mathcal{A}'}(\mathcal{A})\).

\item[(iv)] Let \(P\subseteq\operatorname{Ob}\mathcal{A}'\) be such that every object of \(\mathcal{A}'\) embeds into an object of \(P\).
If \(\eta(X)\) is an isomorphism for all \(X\in P\), and \(F,G\) are way-out right,
then \(\eta(X)\) is an isomorphism for all \(X\in\operatorname{Ob}\mathcal{A}'\).
\end{enumerate}
\end{proposition}

\begin{definition}
Let \(X^\bullet\) be a complex over abelian category \(\mathcal{A}\), \(n\in\mathbb{Z}\).
Define truncation functors:
\begin{align*}
\tau_{>n}(X^\bullet)&\colon \cdots\to 0\to X^{n+1}\to X^{n+2}\to\cdots,\\
\tau_{\le n}(X^\bullet)&\colon \cdots\to X^{n-1}\to X^n\to 0\to\cdots,\\
\sigma_{>n}(X^\bullet)&\colon \cdots\to 0\to \operatorname{im}d^n\to X^{n+1}\to X^{n+2}\to\cdots,\\
\sigma_{\le n}(X^\bullet)&\colon \cdots\to X^{n-1}\to \ker d^n\to 0\to\cdots.
\end{align*}

There are natural morphisms inducing short exact sequences of complexes:
\begin{align}
0 &\to \tau_{>n}(X^\bullet) \to X^\bullet \to \tau_{\le n}(X^\bullet) \to 0,\\
0 &\to \sigma_{\le n}(X^\bullet) \to X^\bullet \to \sigma_{>n}(X^\bullet) \to 0.
\end{align}

The \(\sigma\)-truncations satisfy:
\[
H^i\big(\sigma_{>n}(X^\bullet)\big)
=
\begin{cases}
H^i(X^\bullet) & i>n,\\
0 & i\le n;
\end{cases}
\qquad
H^i\big(\sigma_{\le n}(X^\bullet)\big)
=
\begin{cases}
H^i(X^\bullet) & i\le n,\\
0 & i>n.
\end{cases}
\]
\end{definition}

\begin{lemma}\label{lem:1.7.2}
For any complex \(X^\bullet\) over \(\mathcal{A}\), there exist triangles in \(\DD(\mathcal{A})\):
\begin{align}
\tau_{>n}(X^\bullet) &\to \tau_{\ge n}(X^\bullet) \to X^n \to T\tau_{>n}(X^\bullet),\\
H^n(X^\bullet) &\to \sigma_{\ge n}(X^\bullet) \to \sigma_{>n}(X^\bullet) \to T H^n(X^\bullet).
\end{align}
\end{lemma}

\begin{proof}
Triangle (3) comes from the short exact sequence
\[
0\to\tau_{>n}(X^\bullet)\to\tau_{\ge n}(X^\bullet)\to X^n\to 0.
\]

For (4), define
\(\sigma_{\ge n}'(X^\bullet)\) by
\[
\cdots\to 0\to X^n/\operatorname{im}d^{n-1}\to X^{n+1}\to\cdots.
\]
There is a quasi-isomorphism \(\sigma_{\ge n}(X^\bullet)\to\sigma_{\ge n}'(X^\bullet)\),
and an exact sequence
\[
0\to H^n(X^\bullet)\to\sigma_{\ge n}'(X^\bullet)\to\sigma_{>n}(X^\bullet)\to 0.
\]
This induces the required triangle in \(\DD(\mathcal{A})\).
\end{proof}

\begin{proof}[Proof of Proposition 1.7.1]
(i) For \(X^\bullet\in\DD_{\mathcal{A}'}^b(\mathcal{A})\),
use descending induction on \(n\) to show
\(\eta\big(\sigma_{>n}(X^\bullet)\big)\) is an isomorphism.
For large \(n\), \(\sigma_{>n}(X^\bullet)\) is acyclic, hence trivial.
The induction step follows from triangle (4) and Proposition 1.1(c).
For small \(n\), \(X^\bullet\cong\sigma_{>n}(X^\bullet)\) in \(\DD(\mathcal{A})\), done.

(ii) Let \(X^\bullet\in\DD_{\mathcal{A}'}^+(\mathcal{A})\).
It suffices to show \(H^j(\eta(X^\bullet))\) is an isomorphism for all \(j\).
Given \(j\), pick \(n_1\ge j+2\), choose \(n_2\) from the way-out property for \(F,G\).
Since \(H^i(\sigma_{>n_2}(X^\bullet))=0\) for \(i\le n_2\),
we have \(H^i(F(\sigma_{>n_2}(X^\bullet)))=H^i(G(\sigma_{>n_2}(X^\bullet)))=0\) for \(i<n_1\),
in particular for \(i=j,j+1\).

By the long exact cohomology sequence:
\begin{align*}
H^j\big(F(\sigma_{\le n_2}(X^\bullet))\big) &\stackrel{\sim}{\to} H^j\big(F(X^\bullet)\big),\\
H^j\big(G(\sigma_{\le n_2}(X^\bullet))\big) &\stackrel{\sim}{\to} H^j\big(G(X^\bullet)\big).
\end{align*}
But \(\sigma_{\le n_2}(X^\bullet)\in\DD_{\mathcal{A}'}^b(\mathcal{A})\), so \(\eta\) is an isomorphism there.
Hence \(H^j(\eta(X^\bullet))\) is an isomorphism.

(iii) Decompose \(X^\bullet\) via \(\sigma_{\le0},\sigma_{>0}\), apply (ii) to each piece, then glue via the triangle from (2).

(iv) Use Lemma 1.4.6 to resolve any \(X\in\mathcal{A}'\) by a complex \(I^\bullet\) with terms in \(P\).
Apply (i),(ii) with \(\tau\)-truncations instead of \(\sigma\)-truncations.
\end{proof}

\begin{proposition}\label{prop:1.7.2}
Let \(\mathcal{A}'\subseteq\mathcal{A}\), \(\mathcal{B}'\subseteq\mathcal{B}\) be thick abelian categories,
\(F\colon\DD_{\mathcal{A}'}^+(\mathcal{A})\to\DD(\mathcal{B})\) a triangulated functor.
\begin{enumerate}
\item[(i)] If \(F(X)\in\DD_{\mathcal{B}'}(\mathcal{B})\) for all \(X\in\operatorname{Ob}\mathcal{A}'\),
then \(F(X^\bullet)\in\DD_{\mathcal{B}'}(\mathcal{B})\) for all \(X^\bullet\in\DD_{\mathcal{A}'}^b(\mathcal{A})\).

\item[(ii)] If \(F(X)\in\DD_{\mathcal{B}'}(\mathcal{B})\) for all \(X\in\operatorname{Ob}\mathcal{A}'\)
and \(F\) is way-out right,
then \(F(X^\bullet)\in\DD_{\mathcal{B}'}(\mathcal{B})\) for all \(X^\bullet\in\DD_{\mathcal{A}'}^+(\mathcal{A})\).

\item[(iii)] If \(F\) is way-out both directions, extend to all \(X^\bullet\in\DD_{\mathcal{A}'}(\mathcal{A})\).

\item[(iv)] Let \(P\subseteq\operatorname{Ob}\mathcal{A}'\) such that every object of \(\mathcal{A}'\) embeds into \(P\).
If \(F(X)\in\DD_{\mathcal{B}'}(\mathcal{B})\) for all \(X\in P\) and \(F\) is way-out right,
then \(F(X)\in\DD_{\mathcal{B}'}(\mathcal{B})\) for all \(X\in\operatorname{Ob}\mathcal{A}'\).
\end{enumerate}
\end{proposition}

\begin{proof}
Analogous to Proposition 1.7.1.
\end{proof}

\begin{proposition}\label{prop:1.7.4}
Let \(\mathcal{A}\) have enough injectives, \(F\colon\mathcal{A}\to\mathcal{B}\) additive of cohomological dimension \(\le n\).
Let \(P\subseteq\operatorname{Ob}\mathcal{A}\) be all \(X\) with \(R^iF(X)=0\) for \(i\ne n\).
Assume every object of \(\mathcal{A}\) is a quotient of an object in \(P\).
Set \(G=R^nF\). Then \(\mathbf{R}F,\mathbf{L}G\) exist, and there is a functorial isomorphism
\[
\psi\colon \mathbf{R}F \stackrel{\sim}{\to} \mathbf{L}G[-n],
\]
where \([-n]\) denotes shifting right by \(n\) degrees.
\end{proposition}

\begin{proof}
Existence of \(\mathbf{R}F\) follows from Corollary 1.5.3.
Existence of \(\mathbf{L}G\) follows by dualizing Theorem 1.5.1 for left derived functors.

The class \(P\) satisfies the dual conditions of Lemma 1.4.6 and Corollary 1.5.3.
Let \(\LL\subseteq\KK(\mathcal{A})\) be complexes of objects in \(P\);
the hypotheses of Theorem 1.5.1 hold, so \(\mathbf{L}G\) exists.

Since \(\LL_{\operatorname{Qis}}\to\DD(\mathcal{A})\) is an equivalence,
it suffices to define \(\psi\) on complexes in \(\LL\):
for any \(X^\bullet\in\LL\),
\[
\psi(X^\bullet)\colon \mathbf{R}F(X^\bullet)\to \mathbf{L}G(X^\bullet)[-n],
\]
compatible with morphisms of complexes.

Given \(X^\bullet\in\LL\), choose a Cartan--Eilenberg resolution \(C^{\cdot\cdot}\).
Define the truncated double complex \(C'^{\cdot\cdot}=\sigma_{\le n}^\mathrm{II}(C^{\cdot\cdot})\) by
\[
C'^{p,q}=
\begin{cases}
C^{p,q} & q<n,\\
\ker d_2^{p,n} & q=n,\\
0 & q>n.
\end{cases}
\]

For each \(p\in\mathbb{Z}\), we have an exact sequence
\[
0\to X^p\to C'^{p,0}\to C'^{p,1}\to\cdots\to C'^{p,n-1}\to C'^{p,n}\to 0.
\]
The terms \(C'^{p,0},\dots,C'^{p,n-1}\) are injective; \(X^p\in P\) implies \(C'^{p,n}\) is \(F\)-acyclic.
This resolution computes \(\mathbf{R}F(X^p)\), giving an isomorphism
\[
\varphi(t^p)\colon G(X^p)=R^nF(X^p) \to F(C'^{p,n})\big/\operatorname{im}F(C'^{p,n-1}),
\]
where \(t^p\colon X^p\to C^{p,\bullet}\) is the resolution map.
This induces a morphism \(\alpha^p\colon F(C^{p,n})\to G(X^p)\), and hence a double complex morphism
\[
\alpha\colon F(C'^{\cdot\cdot})\to G(X^\bullet),
\]
with \(G(X^\bullet)\) concentrated in the \(n\)-th row.

Passing to associated simple complexes:
\[
s(\alpha)\colon F\big(s(C'^{\cdot\cdot})\big)\to G(X^\bullet)[-n].
\]

By Lemma 1.7.5(f) below, the augmentation \(u\colon X^\bullet\to s(C^{\cdot\cdot})\) is a quasi-isomorphism into an \(F\)-acyclic complex, yielding an isomorphism
\[
\varphi(u)\colon \mathbf{R}F(X^\bullet)\to F\big(s(C^{\cdot\cdot})\big).
\]
Since \(X^\bullet\) consists of \(G\)-acyclic objects, we have
\[
\varphi(\operatorname{id}_{X^\bullet})\colon G(X^\bullet)\to \mathbf{L}G(X^\bullet).
\]
Composing these gives the desired morphism
\[
\psi(X^\bullet)\colon \mathbf{R}F(X^\bullet)\to \mathbf{L}G(X^\bullet)[-n].
\]

The morphism \(\psi(X^\bullet)\) is independent of the Cartan--Eilenberg resolution choice.
To see \(\psi\) is an isomorphism: \(\mathbf{R}F,\mathbf{L}G\) are way-out in both directions. By the way-out functor lemma (Proposition 1.7.1), reduce to checking \(\psi(X)\) is an isomorphism for \(X\in P\), which holds by construction.
\end{proof}

\begin{proposition}\label{prop:1.7.6}
Let \(\mathcal{A}\) have enough injectives, \(X^\bullet\in\KK^+(\mathcal{A})\). The following are equivalent:
\begin{enumerate}
\item[(i)] \(X^\bullet\) admits a quasi-isomorphism into a bounded complex of injective objects of \(\mathcal{A}\).
\item[(ii)] The functor \(\mathbf{R}\mathcal{H}\textit{om}(\,\cdot\,,X^\bullet)\colon\DD(\mathcal{A})^\circ\to\DD(\mathbf{Ab})\) is way-out left (hence way-out both ways).
\item[(iii)] There exists \(n_0\in\mathbb{Z}\) such that \(\operatorname{Ext}^i(Y,X^\bullet)=0\) for all \(Y\in\operatorname{Ob}\mathcal{A}\) and all \(i>n_0\).
\end{enumerate}
\end{proposition}

\begin{proof}
\((i)\Rightarrow(ii)\): Assume \(X^\bullet\) is bounded injective. Then \(\mathbf{R}\mathcal{H}\textit{om}(Y^\bullet,X^\bullet)\cong\mathcal{H}\textit{om}(Y^\bullet,X^\bullet)\), which is clearly way-out left.

\((ii)\Rightarrow(iii)\): Since \(F=\mathbf{R}\mathcal{H}\textit{om}(\,\cdot\,,X^\bullet)\) is way-out left, choose \(n_0\) such that \(H^i(F(Y^\bullet))=0\) whenever \(H^i(Y^\bullet)=0\) for \(i<n_0\).
For fixed \(Y\in\mathcal{A}\), set \(Y'^\bullet = T^{-n_0}(Y)\). Then \(H^i(Y'^\bullet)=0\) for \(i<n_0\), so \(H^i(F(Y'^\bullet))=0\) for \(i>0\). But
\[
H^i(F(Y'^\bullet))=H^i\big(F(T^{-n_0}Y)\big)=H^{i+n_0}(F(Y)),
\]
hence \(\operatorname{Ext}^{i+n_0}(Y,X^\bullet)=0\) for all \(i>0\), i.e., \(\operatorname{Ext}^j(Y,X^\bullet)=0\) for \(j>n_0\).

\((iii)\Rightarrow(i)\): Take an injective resolution \(X'^\bullet\in\KK^+(\mathcal{A})\) of \(X^\bullet\) (Lemma 1.4.6). Suppose \(H^m(X'^\bullet)\neq0\) for some \(m>n_0\). From the exact sequence
\[
0\to B^m(X'^\bullet)\to Z^m(X'^\bullet)\to H^m(X'^\bullet)\to 0,
\]
the inclusion is proper. Choose \(Y=Z^m(X'^\bullet)\); then
\[
\operatorname{Hom}(Y,B^m(X'^\bullet))\subsetneq \operatorname{Hom}(Y,Z^m(X'^\bullet)).
\]
But \(\operatorname{Ext}^m(Y,X^\bullet)=0\) by hypothesis, which forces the map to be an isomorphism, a contradiction. Thus \(H^i(X'^\bullet)=0\) for all \(i>n_0\).

For \(n\ge n_0\), the truncation map \(i\colon\sigma_{\le n}(X'^\bullet)\to X'^\bullet\) is a quasi-isomorphism.
For \(n\ge n_0\), \(Z^{n+1}(X'^\bullet)=B^{n+1}(X'^\bullet)\) is injective.
Consider the exact sequence
\[
0\to\sigma_{\le n}(X'^\bullet)\to\tau_{\le n}(X'^\bullet)\to B^{n+1}(X'^\bullet)[-n]\to 0.
\]
We show \(\operatorname{Ext}^1(Y,B^{n+1})=0\) for all \(Y\in\mathcal{A}\). We have
\[
\operatorname{Ext}^1(Y,B^{n+1})
=\operatorname{Ext}^{n+1}\big(Y,B^{n+1}[-n]\big).
\]

From the Ext long exact sequence:
\[
\cdots\to\operatorname{Ext}^{n+1}\big(Y,\tau_{\le n}(X'^\bullet)\big)
\to\operatorname{Ext}^{n+1}\big(Y,B^{n+1}[-n]\big)
\to\operatorname{Ext}^{n+2}\big(Y,\sigma_{\le n}(X'^\bullet)\big)\to\cdots.
\]
The left term vanishes because \(\tau_{\le n}(X'^\bullet)\) is bounded injective; the right term vanishes by hypothesis. Hence \(\operatorname{Ext}^1(Y,B^{n+1})=0\), so \(B^{n+1}\) is injective.
Thus \(\sigma_{\le n}(X'^\bullet)\) is a bounded injective complex quasi-isomorphic to \(X^\bullet\).
\end{proof}

\begin{definition}
If \(X^\bullet\in\KK^+(\mathcal{A})\) satisfies the equivalent conditions of Proposition 1.7.6, we say \(X^\bullet\) has \emph{finite injective dimension}.
\end{definition}

\begin{corollary}\label{cor:1.7.7}
The complexes of finite injective dimension form a triangulated localizing subcategory of \(\KK(\mathcal{A})\).
\end{corollary}

\begin{proof}
Closure under triangles follows from the Ext long exact sequence. Localizing follows from Proposition 1.3.3: any complex quasi-isomorphic to a finite injective dimension complex also has finite injective dimension.
\end{proof}

\begin{definition}
Denote by \(\KK^+(\mathcal{A})_{\mathrm{fid}}\) the full subcategory of \(\KK^+(\mathcal{A})\) consisting of complexes of finite injective dimension. The corresponding derived subcategory is \(\DD^+(\mathcal{A})_{\mathrm{fid}}\).
\end{definition}

\textbf{Remark.}
These subcategories differ from \(\KK_{\mathcal{A}'}(\mathcal{A}),\DD_{\mathcal{A}'}(\mathcal{A})\), as they cannot be characterized solely by conditions on the cohomology objects \(H^i(X^\bullet)\).
